\documentclass[11pt,a4paper]{article}

\usepackage[margin=0.6in, top=0.55in, bottom=0.55in]{geometry}
\usepackage{titlesec}
\usepackage{enumitem}
\usepackage{hyperref}
\usepackage{xcolor}
\usepackage{xeCJK}

\setCJKmainfont{Noto Sans CJK SC}

\hypersetup{colorlinks=true, urlcolor=blue!70!black, linkcolor=blue!70!black}

\titleformat{\section}{\large\bfseries}{}{0em}{}[\titlerule]
\titlespacing{\section}{0pt}{12pt}{6pt}

\pagenumbering{gobble}
\setlist{nosep, leftmargin=1.5em, topsep=3pt, itemsep=2pt}

% Custom commands for consistent formatting
\newcommand{\projectentry}[5]{%
    \noindent\textbf{#1} \hfill #2 \\[2pt]
    #4 \hfill {\small #3} \\[1pt]
    #5\par\vspace{8pt}%
}

\newcommand{\projectentrysingle}[3]{%
    \noindent\textbf{#1} \hfill #2 \\[2pt]
    #3\par\vspace{8pt}%
}

\begin{document}

% Header
\begin{center}
    {\Huge\bfseries 辛 杰} \\[8pt]
    [\href{mailto:cs.xinjie@gmail.com}{email}] \quad
    [\href{https://github.com/ftxj}{github}] \quad
    [\href{https://ftxj.github.io}{homepage}]
\end{center}

\vspace{2pt}

% Experience
\section{工作经历}

\noindent\textbf{NVIDIA} \hfill 上海 \\
\textit{Senior Compute Architect} \hfill 2022 -- Present
\begin{itemize}
    \item 研究下一代 GPU（Hopper、Blackwell、Rubin）上的端到端模型性能自动化优化
    \item 开发深度学习编译器（CuTile、Triton-to-CuTile）及 MLPerf 基准测试优化
    \item 与 PyG、Megatron-LM、OpenAI Triton 团队合作开发编译器基础设施
\end{itemize}

\vspace{4pt}

% Education
\section{教育背景}

\noindent\textbf{华中科技大学} \hfill 武汉 \\
\textit{计算机科学与技术 · 硕士} \hfill 2019 -- 2022

\vspace{6pt}

\noindent\textbf{华中科技大学} \hfill 武汉 \\
\textit{计算机科学与技术 · 学士} \hfill 2015 -- 2019

\vspace{4pt}

% Projects
\section{项目经历}

\projectentrysingle{Agent for Compiler}{2026 -- Present}
{设计专为编译器优化的轻量级 harness 系统。}

\projectentry{Triton-to-CuTile}{2024 -- 2025}{[\href{https://github.com/triton-lang/Triton-to-tile-IR}{code}]}
{OpenAI 孵化项目；连接 Triton 前端与 CuTile 后端的编译器桥接层。}
{使 Triton 用户能够使用 NVIDIA 优化的 tile 编程基础设施。}

\projectentry{CuTile}{2023 -- Present}{[\href{https://github.com/NVIDIA/cutile-python}{code}] [\href{https://developer.nvidia.com/blog/tag/cutile/}{blog}]}
{NVIDIA GPU 的 CUDA Tile 编程标准。}
{设计并实现了深度学习工作负载的优化 pass。}

\projectentry{NvFuser}{2022 -- 2023}{[\href{https://github.com/csarofeen/pytorch/commits/devel/?author=ftxj}{code}] [\href{https://pytorch.org/blog/introducing-nvfuser-a-deep-learning-compiler-for-pytorch/}{blog}]}
{PyTorch JIT 编译器，在 GNN 模型上实现 1.2x--2.5x 加速。}
{为稀疏工作负载启用了 gather/scatter/index\_select 图操作。}

\projectentry{OpenFold2 Training}{2023 -- 2023}{[\href{https://github.com/mlcommons/hpc_results_v3.0/tree/main/NVIDIA/benchmarks/openfold}{code}] [\href{https://arxiv.org/abs/2404.11068}{paper}] [\href{https://developer.nvidia.com/blog/optimizing-openfold-training-for-drug-discovery/}{blog}]}
{MLPerf HPC v3.1 榜首，端到端加速 7.5 倍。}
{优化大规模 GPU 集群上的蛋白质结构预测训练。}

\projectentry{GPT-3 Training}{2023 -- 2023}{[\href{https://github.com/NVIDIA/Megatron-LM}{code}]}
{MLPerf Training v4.0；为 175B 模型贡献了 1.02x 训练加速。}
{优化大语言模型的分布式训练流水线。}

\projectentry{PyTorch Geometric TorchScript}{2022 -- 2022}{[\href{https://github.com/pyg-team/pytorch_geometric/commits/master/?author=ftxj}{code}]}
{为 PyG 社区添加 TorchScript 支持，实现 GNN 模型的 JIT 编译。}
{贡献了图神经网络生产部署的核心基础设施。}

\projectentry{SpMM on GPU -- Champion, MIT/IEEE Graph Challenge}{2021 -- 2021}{[\href{https://github.com/CGCL-codes/Graphchallenge21}{code}] [\href{https://ieeexplore.ieee.org/document/9622791}{paper}]}
{高性能稀疏矩阵乘法；相比 cuSPARSE 最高加速 652 倍。}
{获得 2021 年 MIT/IEEE/Amazon Graph Challenge 冠军。}

\vspace{4pt}

% Current Focus
\section{当前方向}

\noindent 我的研究兴趣已全面转向 AI Agent —— 探索 LLM 如何自主优化编译器、自动化性能调优，并变革传统软件开发流程。

\end{document}
